% Fragment: analytic comparison, unitary identification, and full theories.
% No preamble.  The master document supplies theorem-like environments,
% criticalnote/reviewnote, \artifact, \Rep, \Com, \FPdim, \Vec, and \HH.

\section{From vacuum unitarity to complete unitarity}
\label{sec:analytic-unitarity}

The exact strong-generation certificates concern the actual subVOA
$U\subset B\subset H=V_{E_8}$.  In particular, they establish
$\dim U_1=6$ and $U\ne B$; they do not establish positivity on all
modules or on fusion products.  We now separate those questions.
\emph{Vacuum unitarity} is the positive invariant form inherited by the
PCT-stable subVOA $U$.  \emph{Strong unitarity} means that all simple
ordinary $U$-modules admit compatible positive invariant forms.
\emph{Complete unitarity} additionally includes positivity of the
canonical Hermitian fusion forms, giving the VOA representation category
its unitary modular tensor structure.  An abstract ribbon equivalence to
a unitary category does not, by itself, prove this last positivity statement.

We use the preceding algebraic analysis as input: regularity of $U$,
the actual restrictions
\begin{equation}
 B\mathord\downarrow_U=U\oplus X,
 \qquad
 H\mathord\downarrow_U=U\oplus X\oplus Y^{\oplus2},
 \qquad h_X=h_Y=1,
 \label{eq:analytic-branching}
\end{equation}
the twelve simple algebraic $U$-modules, and the
$\operatorname{HI}(\mathbb Z/3)$ fusion ring of the all-module category
for the original extension.  Every use of these inputs inherits the
qualifications attached to their preceding proofs.  Agreement of modular
data alone is never used as a classification of tensor categories.

\begin{criticalnote}{Analytic proof boundary}
We distinguish two comparisons.  The linear reconstruction of modules
and the independent sector count are proved below from the stated
algebraic inputs and external reconstruction theorems.  These resolve
the earlier obligations A3 and A2, respectively.  The remaining
obligation A1 identifies the particular bounded annulus insertions used
in the positivity theorem.  It is stronger than a linear equivalence
of module categories.  Complete unitarity, the unitary ribbon
identification for the VOA, and the specifically $\HH_3$ VOA extension
retain this explicit hypothesis.  The net category itself is identified
unitarily without A1.  The arguments below are mathematical deductions
for review, not an assertion of independent expert validation.
\end{criticalnote}

\subsection{The subnet and its sector count}
\label{subsec:analytic-net-count}

There are two net constructions in use.  Write $\mathcal N_V$ for the
net generated by the closed smeared fields of a strongly local VOA $V$,
and temporarily write $\mathcal A_V^{\mathrm{bl}}$ for the net generated
by Tener's bounded localized insertions.  For $V=U,B,H$ both exist:
strong locality passes to unitary subalgebras, as do bounded localized
vertex operators.  The conformal vectors are the ambient ones, so the
rotations and central charge agree in each construction.

First work with the bounded-insertion inclusion.  Regularity of $U$
and Tener~\cite[Corollary~4.24]{TenFusion} give
\[
 [\mathcal A_H^{\mathrm{bl}}:\mathcal A_U^{\mathrm{bl}}]<\infty.
\]
The ambient net is the holomorphic lattice net, by
Tener~\cite[Proposition~7.5]{TenFields}.  Hence
Longo~\cite[Theorem~24]{Longo} makes $\mathcal A_U^{\mathrm{bl}}$ completely
rational.  Tener's Lemma~7.6 gives a vacuum net isomorphism
$\mathcal A_U^{\mathrm{bl}}\cong\mathcal N_U$, using the ambient
lattice comparison.  It follows that $\mathcal N_U$ is completely
rational as well.  Only this isomorphism-invariant property is
transferred here: we do not infer equality of the two concrete net
constructions or an identification of their embedded extensions.
In particular, this step does not assume strong unitarity of $U$.

The actual smeared-field inclusion $\mathcal N_U\subset\mathcal N_H$
will be used for the sector calculation.  Its branching and index are
established below, after the elementary domain lemma needed for
reconstruction.

\begin{proposition}[The actual net branching and sector count; A2]
\label{clm:analytic-net-count}
Assume the algebraic branching and rank-six ring deductions of
Section~\ref{sec:branching}, including their certificate hypotheses.
The restriction of the holomorphic $H$ net has precisely the three
inequivalent irreducible $\mathcal N_U$ summands corresponding to
\eqref{eq:analytic-branching}, with multiplicities $1,1,2$.
Its boundary fusion category $\mathcal F_{\mathrm{net}}$ has rank six,
untwisted invertible subcategory $\Vec_{\mathbb Z/3}$, and the
$\operatorname{HI}(\mathbb Z/3)$ fusion ring.  Consequently
\[
 \Rep(\mathcal N_U)\simeq Z(\mathcal F_{\mathrm{net}}),
 \qquad
 \#\operatorname{Irr}\Rep(\mathcal N_U)=12.
\]
\end{proposition}
\begin{proof}
Every field $v\in U$ satisfies polynomial energy bounds on $H$, since
it is a field of the strongly local lattice VOA and the conformal
Hamiltonian is unchanged.  Its restriction to an algebraic unitary
summand satisfies the same bound.  Consequently the finite-energy
spaces are cores for the smeared fields in the respective Hilbert
completions.  Apply Lemma~\ref{lem:analytic-graph-reduction} below to
the orthogonal summand projections.  This proves that the algebraic
decomposition of $H$ is a Hilbert direct-sum decomposition for the
actual represented subnet.  Conversely, a subnet intertwiner commutes
with its local diffeomorphisms and hence with $L_0$.  The equality of
closed fields in reconstruction, or Fourier coefficients of their
matrix elements, then makes its finite-energy restriction a
$U$-module map.  Thus simple summands stay irreducible and inequivalent,
and the multiplicities remain $1,1,2$.

The inclusion is irreducible.  To see the local assertion explicitly,
let $x\in\mathcal N_H(I)\cap\mathcal N_U(I)'$.  Locality makes $x$
commute with $\mathcal N_U(I')$ as well.  Strong additivity of
$\mathcal N_U$ shows that the algebras of $I$ and its complementary
interval generate the global represented subnet: remove the two
endpoints from any test interval and apply strong additivity twice.
Thus $x$ commutes with the global subnet and, in particular, its local
diffeomorphisms and common rotations.  Its action on the vacuum lies
in the one-dimensional zero-energy space of $H$, so the separating
property of the vacuum makes it scalar.  The precise finite-index
input is now
Kawahigashi--Longo~\cite[Proposition~2.3]{KL}: an irreducible local extension
of a net with finitely many irreducible DHR sectors, and in which every
sector is a direct sum of finite-statistical-dimension sectors, has
finite index.  These hypotheses follow
from the complete rationality already established for $\mathcal N_U$.
Thus the actual inclusion has finite index.  Its dual canonical sector
is the restriction of the $H$ vacuum computed above, and its dimension
is the inclusion index.  Its algebra $A_{\mathrm{net}}$ is Lagrangian because
$\mathcal N_H$ is holomorphic.  The preceding comparison also gives
the algebra identity
\begin{equation}
 \operatorname{End}(A_{\mathrm{net}})
 \cong\mathbb C\oplus\mathbb C\oplus M_2(\mathbb C).
 \label{eq:analytic-end-algebra}
\end{equation}
Here the multiplication is composition, rather than a dimension count
alone.  Center induction identifies its dimension with
$\operatorname{rank}(\mathcal F_{\mathrm{net}})$, and identifies the
algebra with the complex fusion algebra of
$\mathcal F_{\mathrm{net}}$.  Thus the latter has rank six and is
noncommutative.

We next identify the invertible objects using only the actual
$\mathbb Z/3$ lattice-extension automorphisms of $H/B$.  They act
on the net inclusion and hence on its canonical algebra.  The
automorphism--invertible-object injection of
Lemma~\ref{lem:automorphisms-invertibles} embeds this group in
$\operatorname{Inv}(\mathcal F_{\mathrm{net}})$.  Since the category
has rank six, the latter group has order three or six.  If its order
were six, the category would be pointed of global dimension six.
The intermediate lattice inclusion has $[\mathcal N_H:\mathcal N_B]=3$,
so multiplicativity of index would give
$[\mathcal N_B:\mathcal N_U]=2$.  An irreducible local index-two
extension is the fixed-point extension of its sign involution: its
canonical algebra is $\mathbf1\oplus\beta$, where $\beta^2=\mathbf1$,
and $\mathbf1\oplus(-\mathrm{id}_\beta)$ is its nontrivial algebra
automorphism.  Theorem~6.9 of \cite{CKLW} lifts this vacuum-preserving
net automorphism to a unitary automorphism $g$ of the actual $B$.
Their Proposition~7.6 gives
$\mathcal N_{B^g}=\mathcal N_B^g=\mathcal N_U$, and injectivity of
their unitary-subalgebra/subnet correspondence (Theorem~7.5) gives
$B^g=U$.  This contradicts
Lemma~\ref{lem:no-sign-extension}.  Consequently
\[
 \operatorname{Inv}(\mathcal F_{\mathrm{net}})\cong\mathbb Z/3.
\]
This argument needs no assertion about the full pointwise stabilizer
$\operatorname{Aut}(B/U)$.

The lattice net $\mathcal N_B$ has discriminant metric group
$\mathbb Z/3\times\mathbb Z/3$ with
$q(a,b)=(2a^2+b^2)/3\pmod{\mathbb Z}$.  The extension $H/B$ is its
diagonal isotropic line $L_+=\{(a,a)\}$, and
$L_-=\{(a,-a)\}$ is a transverse isotropic line.  The subcategory
$\mathcal T_-$ supported on $L_-$ is Tannakian and is tensor
equivalent to the untwisted $\Vec_{\mathbb Z/3}$.  Let $Q_+$ be the
extension algebra supported on $L_+$.  Free-module induction
$t\mapsto t\otimes Q_+$ is a tensor functor from $\mathcal T_-$ to
the category of all $Q_+$-modules.  For simple $t,t'\in\mathcal T_-$,
its adjunction gives
\[
 \operatorname{Hom}_{Q_+}(t\otimes Q_+,t'\otimes Q_+)
 \cong\operatorname{Hom}(t,t'\otimes Q_+)
 =\begin{cases}\mathbb C,&t=t',\\0,&t\ne t'.\end{cases}
\]
Here transversality $L_+\cap L_-=\{0\}$ gives the last equality.
The target has dimension $9/3=3$, so these three distinct invertible
images exhaust its simple objects.  Induction is therefore a tensor
equivalence.  Forgetting the intermediate $B$-action embeds this
category fully faithfully and tensorially in
$\mathcal F_{\mathrm{net}}$; relative tensor product over the same
$H$ algebra is unchanged.  Thus the entire pointed subcategory is
untwisted.

We may now apply the rank-six ring and indicator argument of
Section~\ref{sec:branching} to $\mathcal F_{\mathrm{net}}$ itself.
Its hypotheses are the fusion-algebra decomposition above, the three
invertibles, and their untwisted associator.  They have just been
proved for this category, and its spherical structure is positive
because it is a $C^*$ fusion category.  The argument gives
\eqref{eq:rank6-ring} with multiplicity one.  In particular, for
$d=(3+\sqrt{13})/2$ and $s=1+d^2$, the center calculation gives dimensions
\begin{equation}
 1,\ d^2,\ \underbrace{s,s,s,s}_{4},\
 \underbrace{3d,3d,3d,3d,3d,3d}_{6}.
 \label{eq:analytic-center-dimensions}
\end{equation}
The same center-induction and indicator calculation gives the last
six twists with exponents
$\{2,5,6,7,8,11\}/13$.  This is the source of the net count
twelve.  Finally the Lagrangian-algebra theorem gives
$\Rep(\mathcal N_U)\simeq Z(\mathcal F_{\mathrm{net}})$, and the
displayed center dimensions give its rank twelve.  No VOA character
transformation matrix was assigned to the net category.
\end{proof}

For comparison, the first paragraph also follows for the
bounded-insertion inclusion directly from
Tener~\cite[Proposition~5.5 and Theorem~5.6]{TenFields}.  The argument above has the
additional advantage that the later automorphism comparison refers
throughout to the original lattice fields.  The purely categorical
deductions used in the proof are recorded in
\artifact{hi_cft_20260923_next_structure_rank46.md},
\artifact{hi_cft_20260923_next_structure_rank6_indicators.md}, and
\artifact{hi_cft_20260923_next_structure_rank6_center.md}.

\subsection{Reverse reconstruction and exhaustion}
\label{subsec:analytic-reconstruction}

Henriques--Tener~\cite[Theorem~9.5 and Corollary~9.7]{HTNet} reconstruct unitary
modules and identify their closed smeared fields.  Remark~9.8 extends
the construction beyond the finite-multiplicity hypothesis.  This
extension matters: finite-dimensional energy eigenspaces must not be
silently assumed for the net representation one is trying to reconstruct.
Their Theorem~8.9 identifies the reconstructed vacuum VOA of
$\mathcal N_U$ with the original $U$, since $U$ is strongly local.

The closure step can be isolated from the VOA theory.  Its formulation
also explains why a formal mode identity alone is insufficient.

\begin{lemma}[Reduction on a common core]
\label{lem:analytic-graph-reduction}
Let $T_0:D\to\mathcal K$ be a densely defined closable operator and
$P$ an orthogonal projection.  Suppose $PD\subset D$ and
$T_0P\xi=PT_0\xi$ for $\xi\in D$.  Then $P$ reduces
$T=\overline{T_0}$ and $T^*$, and it commutes with the von Neumann
algebra generated by $T$.  In particular, this holds simultaneously
for a family of fields satisfying these hypotheses on a common core.
\end{lemma}
\begin{proof}
If $\xi_j\in D$ converges to $\xi\in D(T)$ in graph norm, then
$P\xi_j\to P\xi$ and $T_0P\xi_j=PT_0\xi_j\to PT\xi$.
Thus $P\xi\in D(T)$ and $TP\xi=PT\xi$.  The same assertion holds
for $1-P$.  It follows that $T$ is the direct sum of its closed
restrictions to $P\mathcal K$ and $(1-P)\mathcal K$; taking adjoints
gives the corresponding direct sum for $T^*$.  Its polar partial
isometry and the spectral projections of $|T|$ are therefore block
diagonal.  These generate the von Neumann algebra of $T$, proving the
last assertion.  Notice that no separate assertion that $D$ is a core
for $T^*$ is required.
\end{proof}

\begin{proposition}[Reverse module comparison; A3]
\label{clm:analytic-reconstruction}
For this regular $U$, reconstruction sends each irreducible
$\mathcal N_U$ sector to an irreducible ordinary unitary $U$-module.
It preserves $L_0$ and induces a fully faithful linear $*$-functor on
finite direct sums,
\[
 \mathcal R:\Rep(\mathcal N_U)\longrightarrow\Rep^{u}(U).
\]
The functor preserves conformal Hamiltonians.  No tensorator or
compatibility with braiding is asserted.
\end{proposition}
\begin{proof}
Let $\mathcal K$ be an irreducible net representation and put
\[
 M=\bigcup_{R>0}1_{[0,R]}(L_0)\mathcal K
\]
as in Henriques--Tener, Remark~9.8.  This is a dense admissible weak
$U$-module with a positive invariant Hermitian form.  Theorem~9.5(ii)
and its extension in that remark identify the closed reconstructed
smeared field with $\pi_I(Y_U(v,f))$ for
$\operatorname{supp}(f)\subset I$.  By Definition~9.4, $M$ is a
core for this closed operator.  The proof of Corollary~9.6 applies
unchanged on bounded-energy vectors, so the VOA zero mode of the
conformal vector is precisely the self-adjoint rotation generator on
this domain.

Regularity decomposes the weak module $M$ algebraically into ordinary
simple modules.  Choose a nonzero simple summand $S$.  Its inherited
form is positive and invariant.  Since its $L_0$ is the restriction
of the net Hamiltonian, its energy decomposition is
\[
 S=\bigoplus_{n\geq0}S_{h+n},\qquad
 \dim S_{h+n}<\infty,\quad h\geq0.
\]
The Hilbert closure $\overline S$ is the orthogonal Hilbert sum of
these finite-dimensional spaces.  In particular it reduces $L_0$,
and the projection $P$ onto $\overline S$ satisfies
\begin{equation}
 P1_{[0,R]}(L_0)\mathcal K
   \subset\bigoplus_{h+n\leq R}S_{h+n}\subset S,
 \qquad PM\subset S\subset M.
 \label{eq:analytic-projection-core}
\end{equation}
This establishes the domain invariance before any commutation with
fields is used.  It requires finite energy multiplicities only for the
chosen ordinary summand, not for $\mathcal K$.

For a homogeneous $v\in U$, the unitary invariance identity expresses
the adjoint of a mode $v_{(n)}$, on $M$, as a finite linear combination
of modes of the vectors $L_1^j\theta v$, where $\theta$ is PCT.
Every such mode preserves $S$.  Testing against $s\in S$ therefore
gives
\[
 \langle Pv_{(n)}\xi,s\rangle
 =\langle\xi,v_{(n)}^{\dagger}s\rangle
 =\langle P\xi,v_{(n)}^{\dagger}s\rangle
 =\langle v_{(n)}P\xi,s\rangle.
\]
Both outer vectors lie in $S$, by
\eqref{eq:analytic-projection-core}, whence
$Pv_{(n)}\xi=v_{(n)}P\xi$ for $\xi\in M$.
The smeared fields on $M$ are Hilbert-space vectors, as in the paragraph
preceding Henriques--Tener Definition~9.4.  Taking their energy
components, or equivalently their convergent Fourier sums, now yields
\[
 P Y^M(v,f)\xi=Y^M(v,f)P\xi,\qquad \xi\in M.
\]
There is no assertion here that a smeared field maps $M$ into itself.
Only its value in $\mathcal K$ and the boundedness of $P$ are needed.
Lemma~\ref{lem:analytic-graph-reduction} extends this identity to the
closures and their adjoints.

The local algebras are generated by the closed smeared fields.
Normal representations preserve this generation statement (equivalently,
apply them to the polar decompositions and spectral projections).
Thus $P$ is a net intertwiner.  Irreducibility and $S\ne0$ imply
$P=1$.  Equation~\eqref{eq:analytic-projection-core} then implies
$M=S$.  This proves ordinaryness and simplicity, including the removal
of the infinite multiplicities permitted in Remark~9.8.

Finally consider morphisms.  A net intertwiner commutes with the
canonical diffeomorphism action, because that action is generated by
the local diffeomorphism groups.  It therefore preserves bounded-energy
vectors.  Commutation with the reconstructed fields gives commutation
with their Fourier coefficients, hence a module map.  Conversely, an
isomorphism of simple ordinary unitary modules is a scalar multiple
of a unitary.  To verify this without presupposing boundedness, compare
the two invariant forms on each finite-dimensional energy space.  The
resulting grade-preserving positive operator is an algebraic
$U$-endomorphism, by invariance of the forms.  An eigenvector in the
finite-dimensional lowest space has an eigenspace that is a nonzero
$U$-submodule.  Simplicity makes that eigenspace the whole module, so
the comparison operator is a single positive scalar on all grades.
The isomorphism therefore has a bounded unitary rescaling.  Its
unitary extension intertwines the smeared fields on their common cores
and hence their closures, so it intertwines the nets.  Between
inequivalent simples both morphism spaces vanish.  Matrices of these
maps handle finite direct sums and respect adjoints.  This proves full
faithfulness and the asserted linear $*$-functor.
\end{proof}

The preliminary version of this deduction is recorded in
\artifact{hi_cft_20260923_next_m3_unitarity_scope.md}.

\begin{proposition}[Strong unitarity from sector exhaustion]
\label{prop:analytic-strong-unitarity}
Under the preceding algebraic inputs, including the count of twelve
simple $U$-modules, every simple $U$-module is
unitarizable, and its character is the character of a net sector.
\end{proposition}
\begin{proof}
Full faithfulness makes the twelve net simples pairwise inequivalent
simple ordinary $U$-modules.  There are exactly twelve algebraic simple
modules, so this injection is a bijection.  Each image carries its
reconstructed positive form, and $L_0$ is preserved.  Finite direct sums
give the corresponding statement for ordinary modules.
\end{proof}

This counting step supplies essential surjectivity only as a linear
comparison.  Neither fusion compatibility, equality of categorical
$S$ matrices, nor positivity of VOA fusion forms follows from it.

\subsection{The annulus compatibility claim}
\label{subsec:analytic-annuli}

The next step concerns the specific geometric annulus system of
Tener~\cite[Definition~3.18 and Corollaries~3.20--3.22]{TenFields}.
Let $(B,A)$ be a localized pair for an interval $I$, with incoming
factor $A$ and outgoing factor $B$.  Subscripts distinguish their
actions on the vacuum and on a reconstructed module.  For an insertion
point $z\in\operatorname{int}(B,A)$, put
\[
 X_v(B,z,A)=\overline{B_UY_U(v,z)A_U}
       \in\mathcal A_U^{\mathrm{bl}}(I).
\]
Here the bar denotes the bounded operator represented by the
sesquilinear insertion form.  At a point outside the incoming annulus
one must not silently treat the unsmeared formal field as an everywhere
defined Hilbert-space operator.  The earlier abstract isomorphism
$\mathcal A_U^{\mathrm{bl}}\cong\mathcal N_U$ does not identify these
operators.  The following claim includes the specific vacuum
identification needed to compare them, and then its module version.

\begin{criticalnote}{A1: new annulus compatibility lemma}
\begin{claim}[Annulus compatibility; unproved geometric comparison]
\label{clm:analytic-annulus}
For the actual geometric system above, choose matching paths and
central lifts in the two annulus constructions.  In the vacuum, the
geometric insertion of Henriques--Tener is $X_v(B,z,A)$ and identifies
$\mathcal A_U^{\mathrm{bl}}$ with $\mathcal N_U$ on the original
vacuum Hilbert space.  If $M$ is an ordinary
unitary module reconstructed from a sector $\pi_M$, then
\begin{equation}
 \pi_M\bigl(X_v(B,z,A)\bigr)
   =\overline{B_MY_M(v,z)A_M},
 \qquad v\in U,\quad z\in\operatorname{int}(B,A).
 \label{eq:analytic-annulus-identity}
\end{equation}
The right side exists as a bounded extension of the insertion form.
\end{claim}
The remaining point requiring verification is a comparison of the
specific geometric annulus systems and their central lifts, including
their localizing diffeomorphisms.  The domain-density and continuation
steps are proved in the lemmas below.  They do not require the
inverse-finite-energy domains to be graph cores.  Agreement after an
unspecified vacuum net isomorphism would not prove the displayed
operator identity.
\end{criticalnote}

The unresolved comparison has a concrete formulation.  For a geometric
incoming annulus $X=(\psi,\rho,t,\gamma)$ in Tener's Definition~3.16,
his vacuum operator is
\[
 U(\psi)e^{-tL(\rho)}U(\gamma)^*.
\]
One must select the corresponding centrally extended annulus
$\underline X$, with the same embedded inner boundary and boundary
parametrizations, and identify this expression with its integrated
action.  The time-ordered integration in
Henriques--Tener~\cite[Theorem~5.5 and formula~(5.2)]{HTAnn} gives the candidate:
constant generators integrate to exponentials, and thin factors to the
diffeomorphism action.  The comparison must retain the central scalar
and use the same path lifts.  A lift normalized only by its action on
the vacuum need not have the same action in a sector of fractional
lowest weight.

More precisely, choose the common localizing diffeomorphism
$\widehat\gamma$ of Tener's Corollary~3.20.  The required geometric
statement is that $\widehat\gamma\underline X$ is the corresponding
localized centrally extended annulus, with its specified lift, and
that its vacuum action is the local Virasoro operator
$U(\widehat\gamma)\pi(X)$.  Once this statement is known,
Henriques--Tener Proposition~4.17 transports it to every sector:
both constructions act there by
\[
 U^M(\widehat\gamma)^*
 \pi^M_I\bigl(U(\widehat\gamma)\pi(X)\bigr).
\]
The same localizing lift must be used for the incoming and outgoing
members of a pair.  This is the precise normalization problem retained
in A1; scalar equality in just one representation is insufficient.

\begin{lemma}[Dense inverse-energy domains]
\label{lem:analytic-inverse-domain}
Let $M$ be an ordinary unitary module and let $A_M,C_M$ be the actions
of geometric annuli on its Hilbert completion.  Suppose
\begin{equation}
 A_MC_M=\lambda r^{L_0},\qquad 0<r<1,\quad\lambda\ne0.
 \label{eq:analytic-round-factorization}
\end{equation}
Then $A_M^{-1}M^{\mathrm{fin}}=C_M M^{\mathrm{fin}}$ is dense.
The same conclusion holds for the outgoing adjoint factor.
\end{lemma}
\begin{proof}
The geometric annulus actions are injective and have dense range by
Henriques--Tener Lemma~4.15.  Equation
\eqref{eq:analytic-round-factorization} first gives
$C_M M^{\mathrm{fin}}\subset A_M^{-1}M^{\mathrm{fin}}$.
Conversely, if $A_M\xi\in M^{\mathrm{fin}}$, then
\[
 A_M\xi
  =A_MC_M\lambda^{-1}r^{-L_0}A_M\xi,
 \qquad
 \xi=C_M\lambda^{-1}r^{-L_0}A_M\xi.
\]
The last equality uses injectivity of $A_M$, and proves the reverse
inclusion.  Finally $M^{\mathrm{fin}}$ is dense and $C_M$ is bounded
with dense range; hence $C_M M^{\mathrm{fin}}$ is dense.  This is the
argument of Henriques--Tener Lemma~7.9 with a central scalar retained.
\end{proof}

For the geometric annuli under consideration the factorization itself
is elementary.  The incoming annulus is embedded in the unit disc,
its smooth inner Jordan domain contains zero, and its outer
parametrization is the identity.  Choose a sufficiently small round
disc $r\mathbb D$ inside that inner domain.  The intervening annulus
$C$ has its outer parametrization chosen to sew to the inner boundary
of $A$, and its inner parametrization is the standard round one.
Then $AC$ is the round annulus.  Choosing compatible lifts gives
\eqref{eq:analytic-round-factorization}.  This proves density after
the annulus-action identification specified above; density is not an
additional unknown analytic assumption.

Define
\begin{equation}
 \mathcal D_A=\{\xi\in\mathcal K_M:A_M\xi\in M^{\mathrm{fin}}\},
 \qquad
 \mathcal D_{B^*}=\{\eta\in\mathcal K_M:B_M^*\eta\in M^{\mathrm{fin}}\}.
\label{eq:analytic-annulus-domains}
\end{equation}
Both domains are dense by Lemma~\ref{lem:analytic-inverse-domain}.

\begin{lemma}[Continuation of the insertion form]
\label{lem:analytic-insertion-continuation}
Let $O_{\mathrm{in}}$ and $O_{\mathrm{out}}$ be the incoming and
outgoing insertion interiors, and put
\[
 \Omega=\operatorname{int}
  (\overline{O_{\mathrm{in}}}\cup\overline{O_{\mathrm{out}}}).
\]
Suppose a bounded geometric insertion $G_v(z)$ is holomorphic on
$\Omega$.  If its matrix elements equal the algebraic insertion form
on $O_{\mathrm{in}}\cup O_{\mathrm{out}}$, for vectors in the two
domains \eqref{eq:analytic-annulus-domains}, the equality holds
throughout $\Omega$.  It uniquely defines a bounded extension of that
form.
\end{lemma}
\begin{proof}
For $\xi\in\mathcal D_A$ and $\eta\in\mathcal D_{B^*}$, consider
\begin{equation}
 \left\langle Y_M(v,z)A_M\xi,B_M^*\eta\right\rangle
 \label{eq:analytic-laurent-pairing}
\end{equation}
as a formal matrix coefficient.  It is a finite Laurent polynomial:
if $v$ has weight $k$, and the input and output energies are $p,q$,
only the mode with $q=p+k-n-1$ can contribute.  Both vectors have
finite energy support.  Thus the form and the geometric matrix
coefficient are holomorphic on $\Omega\subset\mathbb C^\times$.

Every component of $\Omega$ meets
$O_{\mathrm{in}}\cup O_{\mathrm{out}}$.  Indeed, its complement
inside $\Omega$ is contained in the union of the two boundaries;
boundaries of open sets are closed with empty interior, and a finite
union of such sets contains no nonempty open component.  The identity
theorem therefore extends the equality across every component,
including any part of the sewing circle lying in $\Omega$.
For fixed $z$, the form is now the restriction of the bounded operator
$G_v(z)$ to two dense test domains.  This proves existence and
uniqueness of its bounded extension.  A graph-core assertion is neither
used nor needed.
\end{proof}

After the geometric normalization in A1, Henriques--Tener
Theorem~9.5(i) supplies the incoming equality.  Sewing is the
multiplicative action constructed in their Section~6.2; the
$*$-compatibility of Proposition~6.19 supplies the outgoing equality.
Lemma~\ref{lem:analytic-insertion-continuation} then proves the full
identity \eqref{eq:analytic-annulus-identity}.  Thus the operator-domain
and continuation parts of the earlier proposed proof have been
resolved.  The remaining issue is the specified comparison of the
two geometric operator constructions, not an implication from strong
unitarity alone.

\begin{theorem}[Conditional complete unitarity]
\label{thm:analytic-complete-unitarity}
Assume Claim~\ref{clm:analytic-annulus}, together with the preceding
algebraic inputs.  Then $U$ is completely
unitary.
\end{theorem}
\begin{proof}
Proposition~\ref{prop:analytic-strong-unitarity} unitarizes every simple
module.  Claim~\ref{clm:analytic-annulus} supplies the representation
required by
Tener~\cite[Definition~2.18]{TenFusion}, for each such module; finite direct sums
preserve the property.  His Corollary~4.6 then supplies positivity of
the canonical fusion forms.  Strong unitarity must be established first:
the unitarizable-fusion-product hypothesis in Theorem~4.4 cannot be used
to bypass the exhaustion argument.
\end{proof}

Recent comparison results leave this specific hypothesis in place.
Gui's extension comparison assumes Condition~II, including complete
unitarity and energy bounds and strong intertwining for a generating
family of charged fields; his inheritance theorem concerns regular
cosets \cite[Definition~4.15 and Theorems~4.20,~6.11]{GuiCompare}.
It does not give inheritance by the present conformal subalgebra:
the commutant of $U$ in $H$ is the vacuum VOA and its double
commutant is $H$.  Henriques--Tener construct bounded vacuum
multidisc operators and prove the vacuum trace-class condition
\cite[Theorems~3.6 and~4.1]{HTGenusZero}; these statements do not
identify the module insertion operators in
\eqref{eq:analytic-annulus-identity}.  A sufficient remaining target
is to match Tener's incoming and outgoing geometric annuli with the
integrated annuli, retaining their embedded boundaries, boundary
parametrizations, common localizing diffeomorphism, and central path
lifts, and then to identify their vacuum point insertions.  Sector
transport, the dense inverse-energy domains, and continuation proved
above would then give A1.  A linear equivalence of unitary module
categories alone supplies neither this comparison nor the positivity
of the canonical VOA fusion forms.

\section{Unitary categorical identifications}
\label{sec:analytic-category}

\begin{lemma}[All modules over a commutative Q-system]
\label{lem:analytic-unitary-all-modules}
Let $\mathcal C$ be a unitary braided fusion category and $A$ a
connected commutative Q-system in $\mathcal C$.  The category
$\mathcal C_A$ of all right $A$-modules, with relative tensor product,
is a unitarizable fusion category.  This assertion does not require
the modules to be local.
\end{lemma}
\begin{proof}
If $r:M\otimes A\to M$ is a right action, the braiding gives a left
action $l=r\circ c_{A,M}$.  Commutativity of $A$ and the hexagon
identity make these commuting actions, and hence an $A$-bimodule.
Naturality of the braiding makes every right-module morphism a
bimodule morphism for these choices.  This realizes $\mathcal C_A$
as the full subcategory of bimodules with $l=r\circ c_{A,M}$.
The relative tensor product is obtained by splitting the separability
idempotent on $M\otimes N$.  Its induced left and right actions again
satisfy this relation, by the same hexagon identity and the balancing
relation.  Thus the subcategory is closed under relative tensor
products, sums and summands, and its unit is $A$.

Every such bimodule is unitarizable by
Gui~\cite[Theorem~2.28]{Gui}; an isomorphism transporting the actions preserves
the displayed relation.  The unitary relative tensor product and its
associativity and unit constraints are supplied by Gui, Theorem~3.13.
The above full subcategory therefore inherits a $C^*$ tensor structure.
Separability gives semisimplicity, connectedness makes its unit simple,
and the usual module-category construction over a connected separable
algebra gives rigidity and finiteness.  It is consequently a unitary
fusion category up to tensor equivalence.

For clarity, locality is the additional monodromy condition
$r\circ c_{A,M}\circ c_{M,A}=r$.  We have not imposed it.  Only the
local subcategory inherits the ambient braiding by descent.
\end{proof}

\begin{theorem}[Unitary Haagerup-center identification, conditional on A1]
\label{thm:analytic-category}
Under the hypotheses of Theorem~\ref{thm:analytic-complete-unitarity},
and the preceding determination of the all-module fusion ring and its
untwisted pointed subcategory, there is a unitary ribbon equivalence
\[
 \Rep(U)\simeq Z(\HH_3).
\]
For the original embedded extension $H/U$, its all-module category is
tensor equivalent to one of $\HH_2$ and $\HH_3$.
\end{theorem}
\begin{proof}
Put $\mathcal C=\Rep(U)$ and let $A$ be the algebra of the actual
unitary extension $H$.  Complete unitarity allows application of
Gui~\cite[Theorems~2.21 and~2.28]{Gui}: $A$ is a standard commutative
Q-system and its modules are unitarizable.
Lemma~\ref{lem:analytic-unitary-all-modules} makes
$\mathcal F=\mathcal C_A$ a $C^*$ fusion category.  Here $\mathcal C_A$ means
\emph{all} modules.  The local modules form $\mathcal C_A^0\simeq\Vec$
because the extension is holomorphic, so substituting local modules
would lose the Haagerup category.

Its already determined fusion ring is
$\operatorname{HI}(\mathbb Z/3)$.  The actual untwisted pointed
subcategory makes the obstruction $c^{(0,3)}$ vanish, while the other
obstruction is in
\[
 H^1\!\left(\mathbb Z/2,H^2(\mathbb Z/3,\mathbb C^\times)\right)=0.
\]
Indeed, the cyclic-group calculation gives
$H^2(\mathbb Z/3,\mathbb C^\times)=
\mathbb C^\times/(\mathbb C^\times)^3=0$;
every nonzero complex number has a cube root.  Also
$(\mathbb Z/3)/2(\mathbb Z/3)=0$, so no choice of distinguished
noninvertible object introduces an additional classification orbit.
These are the hypotheses in
Izumi~\cite[Definition~2.7 and Theorems~5.9 and~9.1]{Izumi}, giving the two
unitary tensor classes $\HH_2,\HH_3$.  No algebra structure on
$1+\rho$ is inserted as an extra assumption.
Grossman--Snyder~\cite[Theorem~1.4]{GS} place these categories in the same
Morita class.  Their centers are therefore braided equivalent.
The positive spherical structures make this a ribbon equivalence, and
the Lagrangian-center theorem gives
\[
 \mathcal C\simeq Z(\mathcal F)\simeq Z(\HH_3).
\]
Reutter~\cite[Theorem~5]{Reutter} provides the unitary braided realization of the
equivalence, preserving the associated positive ribbon structures.
\end{proof}

This uses the published unitary classification as an external input;
its classification equations have not been replayed here.  It avoids
the broader pivotal census, but it does not prove that the natural
reconstruction functor $\mathcal R$ preserves Connes fusion, the VOA
tensor product, or braiding.  That particular comparison is a distinct
statement and is not needed for the theorem just proved conditionally.

\begin{corollary}[Haagerup-center identification for the conformal net]
\label{cor:analytic-net-haagerup}
Under the preceding algebraic inputs, the completely rational
central-charge-eight net $\mathcal N_U$ satisfies
\[
 \Rep(\mathcal N_U)\simeq Z(\HH_3)
\]
as a unitary ribbon category.  Its actual extension to the $E_8$
lattice net has boundary category equivalent to $\HH_2$ or $\HH_3$.
Neither conclusion requires Claim~\ref{clm:analytic-annulus}.
\end{corollary}
\begin{proof}
Proposition~\ref{clm:analytic-net-count} constructs
$\mathcal F_{\mathrm{net}}$ as a unitary fusion category with the
$\operatorname{HI}(\mathbb Z/3)$ fusion ring and an untwisted pointed
subcategory.  The Izumi classification and the vanishing obstruction
groups used in the preceding proof therefore apply directly to this
category.  Thus $\mathcal F_{\mathrm{net}}$ is tensor equivalent to
$\HH_2$ or $\HH_3$.  Their Morita equivalence, the unitary
Lagrangian-center theorem, and uniqueness of the unitary braided
structure give
$\Rep(\mathcal N_U)\simeq Z(\mathcal F_{\mathrm{net}})\simeq Z(\HH_3)$.
No positivity statement about VOA intertwining-operator spaces is
used in this argument.
\end{proof}

\section{Full net theory and a conditional \texorpdfstring{$E_8$}{E8} VOA extension}
\label{sec:analytic-full-theory}

\begin{corollary}[A chosen holomorphic net extension]
\label{cor:analytic-net-extension}
Under the preceding algebraic inputs, there is an irreducible local
conformal-net extension $\mathcal N_U\subset\mathcal E_3$ of central
charge eight.  It is holomorphic and has index
\[
 [\mathcal E_3:\mathcal N_U]
 =\operatorname{Dim}(\HH_3)=3(1+d^2),
 \qquad d=(3+\sqrt{13})/2.
\]
Its all-module boundary category is $\HH_3$, and its algebra-bimodule
category is $\HH_3\boxtimes\HH_3^{\mathrm{op}}$, with $36$ simple
objects.  This construction does not require A1.
\end{corollary}
\begin{proof}
Transport the canonical connected commutative Q-system
$I_{\HH_3}(\mathbf1)\in Z(\HH_3)$ along
Corollary~\ref{cor:analytic-net-haagerup}; call the result
$Q_3\in\Rep(\mathcal N_U)$.  The local-net reconstruction theorem
of Longo--Rehren~\cite[Theorem~4.9]{LRNets} produces the extension.
Connectedness is precisely the multiplicity-one hypothesis for the
identity sector, hence gives irreducibility.  The Q-system has finite
statistical dimension, which equals the displayed index; its
commutativity is precisely the locality condition in that theorem.
The conformal-net form of the correspondence is recalled in
\cite[Section~2.6]{BisTY}.  The represented local diffeomorphism
unitaries of the subnet give the conformal action of the extension,
so the central charge remains eight.

The finite-index theorem \cite[Theorem~24]{Longo} gives complete
rationality.  Since $Q_3$ is Lagrangian, the index relation gives
\[
 \mu(\mathcal E_3)
 =\frac{\mu(\mathcal N_U)}{\dim(Q_3)^2}=1,
\]
and the extension is holomorphic.  The canonical-algebra equivalences
of \cite[Corollary~4.1 and Proposition~4.4]{DMNO} identify the
all-module category with $\HH_3$.  The corresponding bimodule
category is $\HH_3\boxtimes\HH_3^{\mathrm{op}}$ by the
double-centralizer description \cite[Sections~4 and~8]{Mueger}; its
simple objects are the $6\cdot6$ ordered pairs.  As below, these
algebra-bimodule defects preserve the original $\mathcal N_U$ subnet.
\end{proof}

No identification of $\mathcal E_3$ with the $E_8$ lattice net is made
here.  Holomorphicity and central charge eight are the conclusions
of the cited net theorems; the VOA classification used later does not
by itself classify arbitrary holomorphic conformal nets.  Nor does
the present construction identify a reconstructed VOA extension with
the selected tensor category of VOA modules.

\begin{proposition}[Diagonal full theory]
\label{prop:analytic-diagonal}
Under the preceding algebraic inputs, the completely rational net
$\mathcal N_U$ admits its canonical unitary local diagonal full theory.
With the twelve-module character determination, its partition function is
\begin{equation}
 Z_{\mathrm{diag}}(\tau,\bar\tau)
   =\sum_{i=0}^{11}\chi_i(\tau)\overline{\chi_i(\tau)},
 \qquad c_L=c_R=8.
 \label{eq:analytic-diagonal-partition}
\end{equation}
\end{proposition}
\begin{proof}
Apply the diagonal Longo--Rehren construction, in the identity-coupling
case of
Rehren~\cite[Theorem~1.4, Proposition~1.5, and Corollary~1.6]{Rehren}.
Propositions~\ref{clm:analytic-net-count} and~\ref{clm:analytic-reconstruction}
identify all sector Hilbert spaces with the twelve ordinary
modules and preserve their Hamiltonians.  Hence their traces give
\eqref{eq:analytic-diagonal-partition}.  The annulus compatibility claim
is not needed for this net construction or this linear character matching.
\end{proof}

\begin{theorem}[Chosen $\HH_3$ extension and defects; conditional on A1]
\label{thm:analytic-e8-extension}
Assume the hypotheses of Theorem~\ref{thm:analytic-category} and the
preceding strict positivity of all nonvacuum lowest conformal weights.
There exists a unitary holomorphic conformal extension $H'\supset U$
isomorphic to $V_{E_8}$ whose all-module category is $\HH_3$.
For its extension algebra $A_3\in\Rep(U)$,
\begin{equation}
 \Rep(U)_{A_3}\simeq\HH_3,
 \qquad
 {}_{A_3}\Rep(U)_{A_3}\simeq\HH_3\boxtimes\HH_3^{\mathrm{op}}.
 \label{eq:analytic-extension-defects}
\end{equation}
The second category describes algebra-bimodule defects preserving the
$U$ chiral subalgebra in the corresponding full extension theory and
contains an $\HH_3$ fusion subcategory.  It has $36$ simple objects,
labelled by ordered pairs of simple objects of $\HH_3$.
\end{theorem}
\begin{proof}
Let $I_{\HH_3}$ be the right adjoint of the forgetful functor
$Z(\HH_3)\to\HH_3$.  Transport the canonical algebra
$I_{\HH_3}(\mathbf1)$ along the unitary ribbon equivalence.  By the
canonical Lagrangian-algebra construction, as in
Davydov--M\"uger--Nikshych--Ostrik~\cite[Corollary~4.1 and
Proposition~4.4]{DMNO}, the resulting $A_3$
is connected, commutative, separable, and twist-free, with
\[
 \FPdim(A_3)^2=\FPdim(\Rep(U)),
 \qquad \Rep(U)_{A_3}^{0}\simeq\Vec,
 \qquad \Rep(U)_{A_3}\simeq\HH_3.
\]
Trivial twist makes its constituent conformal weights integral.
Connectedness and strict positivity away from the vacuum give a
one-dimensional weight-zero space.  The extension correspondence of
Huang--Kirillov--Lepowsky~\cite[Theorems~3.2--3.6]{HKL}, together with Gui's
unitary Q-system theorem, produces a unitary conformal VOA extension
$H'$.  Here the transported algebra is a Q-system: the canonical
algebra in the unitary center has its standard positive Frobenius
structure, and the equivalence is unitary.  Thus Gui's theorem is
applied to a specified positive algebra structure, rather than
to an arbitrary separable algebra without a chosen $*$-structure.

For completeness, the finiteness condition on the extension can be
checked at the level of $C_2$ spaces.  Write
$H'=\bigoplus_i M_i^{\oplus m_i}$ as a finite sum of ordinary
$U$-modules and set
\[
 C_2^U(H')=\operatorname{span}\{u_{(-2)}w:u\in U,\ w\in H'\}.
\]
The module $C_2$-cofiniteness theorem of
Abe--Buhl--Dong~\cite[Proposition~5.2]{ABD} implies
$\dim H'/C_2^U(H')<\infty$.  Since
$C_2^U(H')\subset C_2(H')$, the extension is $C_2$-cofinite.
Complete reducibility for the extension follows from the separable
algebra/module correspondence in Huang--Kirillov--Lepowsky,
Theorems~3.5--3.6.  The unitary invariant form makes the simple
extension self-contragredient.  Together with the already verified
CFT-type grading, these are the strong-rationality hypotheses.
The Lagrangian condition gives
$\Rep(H')\simeq\Rep(U)_{A_3}^0\simeq\Vec$, so it is holomorphic.
Its central charge is eight, so
Dong--Mason~\cite[Theorem~1]{DM} identifies $H'$ with $V_{E_8}$.

For the canonical algebra, the algebra-bimodule duality gives the
second equivalence in \eqref{eq:analytic-extension-defects}; see
M\"uger~\cite[Sections~4 and~8]{Mueger}.  One way to keep track of the opposite
is to view $\HH_3$ as a module category over its center.  The two
commuting actions on it are
\[
 (x\boxtimes y)\cdot m=x\otimes m\otimes y.
\]
Composing two such functors reverses the order of the right factors,
which accounts for ``op'' in the formula.  The double-centralizer
description identifies these endofunctors with the bimodule category
of the canonical algebra.  The equivalence transports with $A_3$.
The first-factor embedding $x\mapsto x\boxtimes\mathbf1$ is fully
faithful and tensor, and supplies the stated $\HH_3$ fusion subcategory.
There are $6\cdot6=36$ simple pairs, with dimensions
$\dim(x)\dim(y)$.  As a consistency check, their total squared
dimension is $\operatorname{Dim}(\HH_3)^2
=\operatorname{Dim}(\Rep(U))$, as required by Morita duality.
\end{proof}

\begin{remark}[Extension choice and preserved symmetry]
\label{rem:analytic-extension-scope}
The selection of the canonical $\HH_3$ algebra proves existence of
$H'$ with the specified all-module category.  It does not identify
$A_3$ with the original embedded $H/U$ algebra; the original boundary
category may be $\HH_2$ or $\HH_3$.  Both holomorphic VOAs being
isomorphic to $V_{E_8}$ does not identify their embeddings of $U$.
Moreover, the defects in \eqref{eq:analytic-extension-defects} preserve
$U$ and need not preserve the enlarged $E_8$ chiral algebra.
The local chiral $H'$-module category is $\Vec$.
\end{remark}

The finite certificates and algebraic arguments supply the stated
representation-theory inputs.  The proved comparisons A2--A3 then give
sector exhaustion, strong unitarity, and a unitary Haagerup-center
identification for the conformal net.  The extra annulus normalization
in A1 is still needed for complete unitarity of the VOA, its specified
unitary ribbon structure, and the chosen $\HH_3$ VOA extension.
In particular, the diagonal net theory and its character interpretation
have a weaker dependency than the chosen unitary VOA extension.
The earlier synthesis
sources are
\artifact{hi_cft_20260923_c8two_final_theorem.md} and
\artifact{hi_cft_20260923_c8two_certificate_audit.md}; their unqualified
analytic conclusion is replaced here by the remaining specific annulus
hypothesis and the proved intermediate reductions.
